% !TEX root = ../../LAFROX_AnalisisFinanciero_2_20261113.tex
\section{Evaluación financiera}
\label{sec:e2-evaluacion}

Con una tasa de descuento de 0,998\,\% mensual, el proyecto tiene un VAN de 23.942,31 UF y una TIR de 7,08\,\% mensual, de modo que conviene ejecutarlo. Las cifras salen de las hojas \texttt{PARAMETROS}, \texttt{INDICADORES} y \texttt{SENSIBILIDAD} del modelo (LafroX, 2026f).

Se evalúa el flujo del proyecto, que son los cobros de hitos y cuotas menos los costos de la sección~\ref{sec:e2-costos-precio}. El crédito, el aporte de los socios y los intereses quedan fuera, ya que su costo está dentro de la tasa de descuento. Los montos son netos, antes de impuestos, en un horizonte de 56 meses.

\subsection{Tasa de descuento}
\label{sec:e2-tasa-descuento}

La tasa de descuento es el costo promedio del dinero que financia el proyecto, con la proporción de 80\,\% de crédito y 20\,\% de capital propio del Formulario E-24 (Distribuidora Puelche S.A., 2026a):
\begin{itemize}
  \item El crédito bancario cuesta 0,9\,\% mensual, equivalente a 11,35\,\% anual.
  \item Los socios exigen 18\,\% anual, que en tasa mensual compuesta es 1,389\,\%. La exigencia supera el costo del banco porque los socios asumen más riesgo, y queda bajo el tope de 20\,\%.
\end{itemize}

La tasa ponderada es 0,8 × 0,9\,\% + 0,2 × 1,389\,\% = 0,998\,\% mensual, equivalente a 12,65\,\% anual. El flujo es mensual, así que cada mes se descuenta con la tasa mensual.

\subsection{Flujo de caja proyectado}
\label{sec:e2-flujo-proyectado}

La Tabla~\ref{tab:e2-flujo-anual} resume el flujo del proyecto por año de contrato. El flujo con financiamiento se ve en la sección~\ref{sec:e2-curva-s}.

{\renewcommand{\LFXfiletefila}{\noalign{{\color{lafroxLinea}\hrule height 0.3pt}\nobreak}}%
\begin{tablalafrox}{Flujo del proyecto por año de contrato, en UF. Las sumas pueden diferir en la última cifra por redondeo. Fuente: modelo financiero, hoja INDICADORES.}{tab:e2-flujo-anual}{Y R{2.8cm} R{2.8cm} R{2.8cm}}
  {\cab{Período} & \cab{Ingresos} & \cab{Egresos} & \cab{Flujo neto}}
  Año 1, meses 1 a 12 & 45.433,67 & 47.893,13 & $-$2.459,46 \\
  Año 2, meses 13 a 24 & 69.730,06 & 52.590,63 & 17.139,43 \\
  Año 3, meses 25 a 36 & 42.600,07 & 36.389,91 & 6.210,16 \\
  Año 4, meses 37 a 48 & 42.600,07 & 34.788,68 & 7.811,39 \\
  Año 5, meses 49 a 56 & 28.400,05 & 22.825,37 & 5.574,68 \\
  \midrule
  \textbf{Total} & \textbf{228.763,92} & \textbf{194.487,71} & \textbf{34.276,21} \\
\end{tablalafrox}}

Solo el primer año cierra en negativo, ya que en él se compra la infraestructura y parten los equipos antes de los primeros hitos. El segundo año deja el mayor flujo, con siete de los doce hitos y el inicio de la operación. Desde el tercer año el flujo depende de la cuota fija y se mantiene estable.

La Tabla~\ref{tab:e2-flujo-mensual} muestra los 56 meses. Los hitos H1 a H11 entran entre los meses 2 y 19, el mes 21 recibe H12 y la primera cuota, y desde el mes 22 entra solo la cuota. El acumulado toca su mínimo en el mes 9 y es positivo desde el mes 17.

\begin{tablalafrox}{Flujo de caja mensual del proyecto, en UF. Las sumas pueden diferir en la última cifra por redondeo. Fuente: modelo financiero, hoja INDICADORES.}{tab:e2-flujo-mensual}{Y R{2.6cm} R{2.6cm} R{2.6cm} R{2.6cm}}
  {\cab{Mes} & \cab{Ingresos} & \cab{Egresos} & \cab{Flujo del mes} & \cab{Acumulado}}
  1 & 0,00 & 2.020,66 & $-$2.020,66 & $-$2.020,66 \\
  2 & 8.077,10 & 15.775,12 & $-$7.698,02 & $-$9.718,68 \\
  3 & 0,00 & 2.416,43 & $-$2.416,43 & $-$12.135,11 \\
  4 & 7.067,46 & 3.111,02 & 3.956,44 & $-$8.178,67 \\
  5 & 0,00 & 5.467,84 & $-$5.467,84 & $-$13.646,52 \\
  6 & 10.096,37 & 4.241,92 & 5.854,45 & $-$7.792,07 \\
  7 & 0,00 & 4.114,60 & $-$4.114,60 & $-$11.906,67 \\
  8 & 0,00 & 2.637,29 & $-$2.637,29 & $-$14.543,96 \\
  9 & 0,00 & 2.421,43 & $-$2.421,43 & $-$16.965,40 \\
  10 & 10.096,37 & 2.264,38 & 7.831,99 & $-$9.133,41 \\
  11 & 0,00 & 2.042,63 & $-$2.042,63 & $-$11.176,04 \\
  12 & 10.096,37 & 1.379,79 & 8.716,58 & $-$2.459,46 \\
  13 & 5.048,19 & 4.749,13 & 299,06 & $-$2.160,40 \\
  14 & 5.048,19 & 4.679,99 & 368,19 & $-$1.792,21 \\
  15 & 0,00 & 4.724,20 & $-$4.724,20 & $-$6.516,41 \\
  16 & 10.096,37 & 5.191,04 & 4.905,33 & $-$1.611,09 \\
  17 & 10.096,37 & 4.428,88 & 5.667,49 & 4.056,41 \\
  18 & 10.096,37 & 3.468,49 & 6.627,88 & 10.684,28 \\
  19 & 5.048,19 & 5.122,63 & $-$74,45 & 10.609,84 \\
  20 & 0,00 & 5.188,57 & $-$5.188,57 & 5.421,27 \\
  21 & 13.646,38 & 4.854,42 & 8.791,96 & 14.213,23 \\
  22 & 3.550,01 & 4.003,94 & $-$453,93 & 13.759,30 \\
  23 & 3.550,01 & 3.164,01 & 385,99 & 14.145,29 \\
  24 & 3.550,01 & 3.015,32 & 534,68 & 14.679,97 \\
  25 & 3.550,01 & 3.446,71 & 103,30 & 14.783,27 \\
  26 & 3.550,01 & 2.982,93 & 567,08 & 15.350,35 \\
  27 & 3.550,01 & 2.926,99 & 623,02 & 15.973,36 \\
  28 & 3.550,01 & 2.974,56 & 575,45 & 16.548,81 \\
  29 & 3.550,01 & 2.923,75 & 626,26 & 17.175,07 \\
  30 & 3.550,01 & 2.952,31 & 597,70 & 17.772,77 \\
  31 & 3.550,01 & 2.960,24 & 589,77 & 18.362,54 \\
  32 & 3.550,01 & 3.321,43 & 228,58 & 18.591,11 \\
  33 & 3.550,01 & 2.941,88 & 608,12 & 19.199,23 \\
  34 & 3.550,01 & 2.833,19 & 716,81 & 19.916,05 \\
  35 & 3.550,01 & 3.266,09 & 283,91 & 20.199,96 \\
  36 & 3.550,01 & 2.859,83 & 690,18 & 20.890,13 \\
  37 & 3.550,01 & 3.303,66 & 246,35 & 21.136,48 \\
  38 & 3.550,01 & 2.811,44 & 738,57 & 21.875,05 \\
  39 & 3.550,01 & 2.834,19 & 715,82 & 22.590,86 \\
  40 & 3.550,01 & 2.841,27 & 708,74 & 23.299,60 \\
  41 & 3.550,01 & 2.778,71 & 771,29 & 24.070,89 \\
  42 & 3.550,01 & 2.835,23 & 714,78 & 24.785,67 \\
  43 & 3.550,01 & 2.853,25 & 696,76 & 25.482,44 \\
  44 & 3.550,01 & 3.200,40 & 349,61 & 25.832,04 \\
  45 & 3.550,01 & 2.767,20 & 782,80 & 26.614,84 \\
  46 & 3.550,01 & 2.680,22 & 869,79 & 27.484,63 \\
  47 & 3.550,01 & 3.148,81 & 401,19 & 27.885,82 \\
  48 & 3.550,01 & 2.734,30 & 815,70 & 28.701,53 \\
  49 & 3.550,01 & 3.456,41 & 93,60 & 28.795,12 \\
  50 & 3.550,01 & 2.688,81 & 861,19 & 29.656,32 \\
  51 & 3.550,01 & 2.676,96 & 873,04 & 30.529,36 \\
  52 & 3.550,01 & 2.723,03 & 826,98 & 31.356,34 \\
  53 & 3.550,01 & 2.629,82 & 920,18 & 32.276,52 \\
  54 & 3.550,01 & 2.747,89 & 802,12 & 33.078,64 \\
  55 & 3.550,01 & 2.719,37 & 830,64 & 33.909,28 \\
  56 & 3.550,01 & 3.183,07 & 366,93 & 34.276,21 \\
  \midrule
  \textbf{Total} & \textbf{228.763,92} & \textbf{194.487,71} & \textbf{34.276,21} & \\
\end{tablalafrox}

\subsection{VAN y TIR}
\label{sec:e2-van-tir}

El VAN es de 23.942,31 UF, o \$957.692.536 y USD~1.064.103. Resulta de restar al valor actual de los ingresos, 176.787,03 UF, el de los sueldos, 107.028,07 UF, y el de los otros egresos, 45.816,65 UF. Al ser positivo, el proyecto gana más de lo que cuesta el dinero que lo financia.

La TIR es de 7,08\,\% mensual, o 127,13\,\% anual. Es alta porque la inversión propia dura poco: el acumulado cae hasta $-$16.965,40 UF en el mes 9 y ya es positivo en el mes 17, porque los hitos pagan casi toda la implementación. Por eso la decisión se basa en el VAN y la TIR es un dato complementario. La TIR tampoco es el margen, que es de 14,98\,\% del precio (sección~\ref{sec:e2-precio-margenes}).

\subsection{Análisis de sensibilidad}
\label{sec:e2-sensibilidad}

Se movió una variable a la vez en $-20$\,\%, $-10$\,\%, $+10$\,\% y $+20$\,\%, dejando las demás en su valor base. La Tabla~\ref{tab:e2-sensibilidad} muestra el VAN de cada caso.

{\renewcommand{\LFXfiletefila}{\noalign{{\color{lafroxLinea}\hrule height 0.3pt}\nobreak}}%
\begin{tablalafrox}{VAN del proyecto ante variaciones de cada variable, en UF. VAN base: 23.942,31 UF. Fuente: modelo financiero, hoja SENSIBILIDAD.}{tab:e2-sensibilidad}{Y R{2.1cm} R{2.1cm} R{2.1cm} R{2.1cm}}
  {\cab{Variable} & \cab{$-$20\,\%} & \cab{$-$10\,\%} & \cab{+10\,\%} & \cab{+20\,\%}}
  Ingresos & $-$11.415,09 & 6.263,61 & 41.621,02 & 59.299,72 \\
  Costo de las HH & 45.347,93 & 34.645,12 & 13.239,51 & 2.536,70 \\
  Otros egresos & 33.105,64 & 28.523,98 & 19.360,65 & 14.778,98 \\
  Tasa de descuento & 25.653,33 & 24.779,43 & 23.139,91 & 22.370,30 \\
\end{tablalafrox}}

El VAN depende sobre todo de los ingresos. Si bajan 10\,\%, el VAN cae 74\,\%, y con una baja de 20\,\% queda negativo. Después viene el costo de las HH: con un alza de 20\,\% el VAN baja a 2.536,70 UF, pero sigue positivo. Los otros egresos y la tasa pesan menos, y con un alza de 20\,\% en cualquiera de ellos el VAN queda sobre 14.700 UF. Los riesgos más relevantes son, por lo tanto, los atrasos en la aprobación de hitos y el aumento de horas, que cubre la reserva de gestión del plan de riesgos (LafroX, 2026e).

\subsection{Punto de equilibrio}
\label{sec:e2-equilibrio}

El punto de equilibrio se calculó en ingresos, en sueldos y en tiempo:
\begin{itemize}
  \item El VAN llega a cero con el 86,46\,\% de los ingresos proyectados. Sin descontar, los costos se cubren con el 85,02\,\% del precio, 194.487,71 UF.
  \item El VAN llega a cero si los sueldos suben 22,37\,\%.
  \item El flujo acumulado es positivo desde el mes 17 y no vuelve a ser negativo. Desde ese mes, lo cobrado supera todo lo gastado.
\end{itemize}
